% ===================== BASES DE DISENO =====================
\section{Bases de Diseno}
\label{sec:bases}

[Este capitulo reune la totalidad de las variables suministradas por el cliente,
los datos normativos, tecnicos y ambientales requeridos para ejecutar los calculos
y analisis del presente estudio. Incluir referencia a documentos contractuales,
especificaciones de equipos y normativas aplicables.]

% ─────────────────────────────────────────────────────────────────────────────
\subsection{Restricciones de Operacion}
\label{ssec:restricciones}

[Definir las restricciones de operacion establecidas por el cliente y los criterios
de aceptacion del analisis.]

\begin{table}[htbp]
    \centering
    \caption{[Leyenda descriptiva de la tabla de restricciones de diseno.]}
    \label{tab:restricciones}
    \setcounter{itemcount}{0}
    \begin{tabularx}{\textwidth}{@{}NXccl@{}}
        \toprule
        \multicolumn{1}{c}{\textbf{Item}} & \textbf{Parametro} & \textbf{Valor Establecido} & \textbf{Valor Recomendado} & \textbf{Referencia} \\
        \midrule
        & <Parametro 1> & <Valor> & <Valor> & <Norma / Cliente> \\
        & <Parametro 2> & <Valor> & <Valor> & <Norma / Cliente> \\
        & <Parametro 3> & <Valor> & <Valor> & <Norma / Cliente> \\
        \bottomrule
    \end{tabularx}
\end{table}

% ─────────────────────────────────────────────────────────────────────────────
\subsection{Condiciones Ambientales del Sitio}
\label{ssec:amb}

[Describir las condiciones ambientales empleadas en los calculos, fundamentadas
en registros historicos de la ubicacion de la planta. Separar condiciones exteriores
de interiores si aplica.]

\begin{table}[htbp]
    \centering
    \caption{[Leyenda descriptiva de condiciones ambientales exteriores.]}
    \label{tab:amb_ext}
    \setcounter{itemcount}{0}
    \begin{tabularx}{\textwidth}{@{}NXccc@{}}
        \toprule
        \multicolumn{1}{c}{\textbf{Item}} & \textbf{Parametro} & \textbf{Minimo} & \textbf{Maximo} & \textbf{Unidad} \\
        \midrule
        & <Parametro 1> & <Valor> & <Valor> & [Unidad] \\
        & <Parametro 2> & <Valor> & <Valor> & [Unidad] \\
        \bottomrule
    \end{tabularx}
\end{table}

% ─────────────────────────────────────────────────────────────────────────────
\subsection{Propiedades Fisicoquimicas del Fluido de Proceso}
\label{ssec:fluido}

[Describir el fluido de proceso y presentar sus propiedades fisicoquimicas a las
temperaturas relevantes de operacion. Citar fuente de los datos (CoolProp, NIST,
DIPPR, handbook con edicion).]

\begin{table}[htbp]
    \centering
    \caption{[Leyenda descriptiva de propiedades del fluido.]}
    \label{tab:propfluido}
    \setcounter{itemcount}{0}
    \begin{tabularx}{\textwidth}{@{}NXcccc@{}}
        \toprule
        \multicolumn{1}{c}{\textbf{Item}} & \textbf{Propiedad} & \textbf{[T1]} & \textbf{[T2]} & \textbf{Unidad} \\
        \midrule
        & [Propiedad 1] & <Valor> & <Valor> & [Unidad] \\
        & [Propiedad 2] & <Valor> & <Valor> & [Unidad] \\
        & [Propiedad 3] & <Valor> & <Valor> & [Unidad] \\
        \bottomrule
    \end{tabularx}
\end{table}

% ─────────────────────────────────────────────────────────────────────────────
\subsection{Especificacion de Tuberias y Accesorios}
\label{ssec:tuberia}

[Describir el sistema de tuberias: materiales, norma de fabricacion, acabados
superficiales, diametros nominales y especificaciones relevantes para los calculos.
Incluir tabla resumen de dimensiones.]

\begin{table}[htbp]
    \centering
    \caption{[Leyenda descriptiva de dimensiones de tuberias.]}
    \label{tab:tuberia}
    \setcounter{itemcount}{0}
    \begin{tabularx}{\textwidth}{@{}NXcccc@{}}
        \toprule
        \multicolumn{1}{c}{\textbf{Item}} & \textbf{Tramo / Funcion} & \textbf{D. Nominal} & \textbf{D. Exterior} & \textbf{D. Interior} & \textbf{Espesor} \\
        & & \textbf{[in]} & \textbf{[mm]} & \textbf{[mm]} & \textbf{[mm]} \\
        \midrule
        & [Tramo 1] & <Valor> & <Valor> & <Valor> & <Valor> \\
        & [Tramo 2] & <Valor> & <Valor> & <Valor> & <Valor> \\
        \bottomrule
    \end{tabularx}
\end{table}

% ─────────────────────────────────────────────────────────────────────────────
\subsection{Especificacion de Equipos Principales}
\label{ssec:equipos}

[Describir los equipos principales del sistema: bombas, intercambiadores, reactores,
tanques, filtros, etc. Para cada equipo incluir: fabricante, modelo, parametros
nominales de operacion, y curvas caracteristicas si aplica.]

\begin{table}[htbp]
    \centering
    \caption{[Leyenda descriptiva del equipo principal.]}
    \label{tab:equipo}
    \setcounter{itemcount}{0}
    \begin{tabularx}{\textwidth}{@{}NXcc@{}}
        \toprule
        \multicolumn{1}{c}{\textbf{Item}} & \textbf{Parametro} & \textbf{Valor} & \textbf{Unidad} \\
        \midrule
        & [Fabricante / Modelo] & <Valor> & [---] \\
        & <Parametro 1> & <Valor> & [Unidad] \\
        & <Parametro 2> & <Valor> & [Unidad] \\
        \bottomrule
    \end{tabularx}
\end{table}

[Incluir figura de curva caracteristica si aplica:]

% \begin{figure}[htbp]
%     \centering
%     \includegraphics[width=0.85\textwidth]{assets/nombre_figura.png}
%     \caption{[Leyenda descriptiva de la figura, incluyendo variables y condiciones.]}
%     \label{fig:equipo}
% \end{figure}

% ─────────────────────────────────────────────────────────────────────────────
\subsection{Normativas y Estandares Aplicables}
\label{ssec:normas}

[Relacionar el marco normativo y estandares tecnicos aplicables al proyecto.]

\begin{table}[htbp]
    \centering
    \caption{[Leyenda descriptiva del marco normativo.]}
    \label{tab:normas}
    \setcounter{itemcount}{0}
    \begin{tabularx}{\textwidth}{@{}NX>{\raggedright\arraybackslash}X>{\raggedright\arraybackslash}Xl@{}}
        \toprule
        \multicolumn{1}{c}{\textbf{Item}} & \textbf{Norma / Estandar} & \textbf{Alcance} & \textbf{Edicion} \\
        \midrule
        & [Norma 1] & [Alcance de aplicacion] & [Edicion / Anio] \\
        & [Norma 2] & [Alcance de aplicacion] & [Edicion / Anio] \\
        & [Norma 3] & [Alcance de aplicacion] & [Edicion / Anio] \\
        \bottomrule
    \end{tabularx}
\end{table}
